\documentclass[11pt]{article}
\usepackage[margin=1in]{geometry}
\usepackage{amsmath,amssymb,amsthm,mathtools}
\usepackage{enumitem}
\usepackage{booktabs}
\usepackage{microtype}
\usepackage{xcolor}
\usepackage{hyperref}
\usepackage{cleveref}
\hypersetup{colorlinks=true, linkcolor=blue, citecolor=blue, urlcolor=blue}

\newtheorem{definition}{Definition}
\newtheorem{theorem}{Theorem}
\newtheorem{proposition}{Proposition}
\newtheorem{lemma}{Lemma}
\theoremstyle{remark}
\newtheorem{remark}{Remark}

\newcommand{\Acc}{\mathsf{ACCEPT}}
\newcommand{\Rej}{\mathsf{REJECT}}
\newcommand{\Unk}{\mathsf{UNKNOWN}}
\newcommand{\Eval}{\mathsf{Eval}}
\newcommand{\Clos}{\mathsf{Clos}}
\newcommand{\Pend}{\mathsf{PENDING}}
\newcommand{\Auth}{\mathsf{Auth}}
\newcommand{\Dis}{\mathsf{Discharge}}
\newcommand{\Combine}{\mathbin{\otimes}}
\newcommand{\Canon}{\mathsf{Canon}}
\newcommand{\Sel}{\mathsf{Sel}}
\newcommand{\Lineage}{\mathsf{Lineage}}
\newcommand{\Supp}{\mathsf{Supp}}
\newcommand{\Sep}{\mathsf{Sep}}

\title{EvidenceVault: Portable Acceptability, Replay Kernel, and Receipt Algebra}
\author{}
\date{March 19, 2026}

\begin{document}
\maketitle

\begin{abstract}
Portable evidence should let a downstream party replay one fixed acceptability judgement
using a carried evidence bundle rather than trust in the original producer. This paper is the series spine and replay kernel for that judgement.
Portable acceptability is a fail-closed replay problem over four closure obligations: \emph{identity closure}
(what artifact is being judged), \emph{semantic closure} (what the judgement means), \emph{selection closure}
(which latent degrees of freedom were exercised), and \emph{transparency closure} (what external consistency evidence is required).
The kernel fixes a compact normal form for portable replay carriers and outcome-bearing decision objects: one replay carrier that keeps the four replay-fixing coordinates together with one carried bridge replay slice, together with a conservative rule for portable plurality via disclosed lineage, dependency-qualified support counting, a conservative outcome interface with applicable-counterevidence and discharge precedence, and a receipt algebra for conjunction, weakening, thresholds, and sequential accumulation once replay and selection surfaces are fixed.
Local carriers may instantiate those rules differently, but portability always reduces to the same replay question on one fixed scope.
Interpretation profiles, authority surfaces, omission commitments, and decision-local replay context are treated throughout as already-carried replay preconditions.
\end{abstract}

\section{Portable acceptability, replay carriers, and scope-bound replay identity}
A governance claim is not yet portable merely because it is signed. The series spine first fixes the \emph{right replay boundary}:
what bytes are being judged, what rules interpret them, which latent choices were fixed or shown irrelevant, and what anti-equivocation evidence must be checked.

\begin{definition}[Portable acceptability instance]
A portable acceptability instance is a tuple
\[
I=(x,\pi,S,T,E),
\]
where $x$ is the artifact or release under judgement, $\pi$ is the policy or verifier profile, $S$ is the set of latent selection variables
(attempt ids, retries, suite choices, randomness knobs, menu selections), $T$ is the transparency context (logs, witnesses, freshness windows,
checkpoints), and $E$ is the carried evidence bundle.
\end{definition}

\begin{definition}[Fail-closed replay judgement]
The replay kernel computes
\[
\Eval(I)\in\{\Acc,\Rej,\Unk\}
\]
by validating the evidence in $E$ against $(x,\pi,S,T)$. Missing or unverifiable prerequisites yield $\Unk$, not $\Acc$.
\end{definition}

\begin{remark}[Standing replay convention]
This paper uses $\Unk$ as the coarse non-authorizing state. Later sections make the conservative outcome interface explicit: local papers may refine that frontier into labels such as $\Pend$ when the blockage is temporal or procedural, but only $\Acc$ ever authorizes portable downstream reuse. Interpretation profiles, omission commitments, temporal guards, checkpoint coherence, issuer rosters, and signer authority are treated here as fixed replay preconditions. Within that fixed replay state, realized lineage, stopping surfaces, attempt-ledger/search-budget accounting, dependency-diverse support counting, and applicable-counterevidence precedence are checked explicitly by the kernel.
\end{remark}

\paragraph{Replay carriers, decision objects, and scope-bound replay identity.} Portable reuse is never authorized by a bare signed bit such as ``pass,'' ``admit,'' or ``accept.''
The common portable unit is first a \emph{replay carrier} that fixes the replay boundary; an outcome-bearing decision object is then a local reading of that same carried replay boundary.

\begin{definition}[Portable replay carrier]
For a portable acceptability instance $I=(x,\pi,S,T,E)$, a portable replay carrier is the tuple
\[
K=((u,\varepsilon,\lambda,\tau),\Sigma),
\]
where $u$ names the judged scope (typically a canonical digest of $x$ together with its decision scope),
$\varepsilon$ names the carried evaluator-semantics coordinate that gives the judgement meaning,
$\lambda$ records either the realized lineage for latent selection variables in $S$ or carried evidence that the judgement is branch-invariant over them,
$\tau$ names the required anchor/history view and validity conditions,
and $\Sigma$ is one carried bridge replay slice of $E$.
\end{definition}

\begin{definition}[Portable decision object]
A portable decision object is an outcome-bearing reading of one portable replay carrier:
\[
D=(K,o),
\]
where $K$ is the carried replay boundary and $o$ is the claimed outcome.
\end{definition}

\begin{proposition}[Kernel normal form for portable replay carriers]
If a local paper's portable object is replay-bearing, then it admits a conservative normalization to a portable replay carrier $K$.
If it additionally claims an authorizing or blocking verdict, then it likewise admits a conservative normalization to a portable decision object $D$ whose replay under the kernel yields the same authorization bit.
Conversely, any safe local extension of $K$ or $D$ may add family-specific diagnostics, batching, or non-authorizing refinements,
but it may not weaken the replay boundary carried by $K$.
\end{proposition}

\begin{proof}
Take $u$, $\varepsilon$, $\lambda$, and $\tau$ from the carried bundle's judged scope, evaluator semantics, selection-closure material, and
transparency context, and let $\Sigma$ be the carried bridge replay slice of evidence.
That already fixes the replay-bearing carrier $K$, with the slice inheriting the scope fixed by those four replay-fixing coordinates.
If the local object also claims an outcome $o$, the kernel then recomputes the authorization bit for that decision-object reading $D$ from the same coordinates.
Extra local fields are safe only if they do not let a bundle authorize, block, or narrate replay after one of those kernel coordinates has been deleted,
weakened, or silently reinterpreted.
That is exactly conservative refinement at the carrier or decision-object level, and it preserves the same replay carrier boundary.
\end{proof}

The judged-scope coordinate $u$ above is intentionally abstract; portable scope binders, namespace pins, and retarget-safe local naming are fixed by the interpretation bridge's portable scope-binding map, while the kernel uses only the resulting judged scope.


\section{Selection closure and portable plurality}
Selection closure is the kernel rule for latent choice over evidentiary paths.
Many EvidenceVault decisions do not depend on one fixed proof alone: a producer may choose a monitor from a menu, retry a failed proving run, escalate from a fast gate to a fuller audit path, stop a sequential process after one prefix, or count several receipts as stronger support only if they are really distinct.
Those are not separate portability theories. They are one kernel question: which admissible path actually ran, and which carried supports really count as different under policy?

\begin{definition}[Selection surface]
A selection surface for an acceptance procedure is a carried tuple
\[
  \Sel=(M,\beta,\varphi,\varrho,\Pi)
\]
where $M$ is the admissible menu of branches, vendors, monitors, judges, proof systems, or retryable actions,
$\beta$ is the budget or cap policy on how often those options may be exercised,
$\varphi$ is the stopping or escalation rule that determines when evidence collection halts and a decision is emitted,
$\varrho$ is the randomness or resampling policy relevant to branch selection or retries,
and $\Pi$ is the interpretation/replay profile under which menu labels and outcomes are read.
A selection surface is portable only when these coordinates are replayable enough for downstream policy to determine which latent choices were allowed and what their semantics were.
\end{definition}

\begin{definition}[Realized lineage and branch invariance]
For a carried selection surface $\Sel$, a realized lineage is replayable evidence
\[
\Lineage(E,\Sel)
\]
that records the actual branch path, retry history, stopping event, and menu selections consumed in producing the supportive bundle $E$.
A branch-invariance certificate is replayable evidence that the policy-relevant outcome is unchanged across a specified set of admissible branch choices under the same $\Sel$.
\end{definition}

\begin{definition}[Attempt ledger and search-budget certificate]
An attempt ledger for $\Sel$ is carried evidence recording the realized attempts, retries, or redraws consumed under that surface together with their scope bindings and outcomes.
A search-budget certificate for a claim under $\Sel$ is carried evidence showing either
(i) a complete attempt ledger for the policy-relevant search surface, or
(ii) a replayable upper bound on undisclosed admissible search under the same decision scope and replay context.
A search-budget certificate is not permission to hide arbitrary search; it is the carried bound that prevents hidden search from being narrated later as if it had always been within policy.
\end{definition}

\begin{definition}[Dependency profile and separation certificate]
A dependency profile for an evidence object $e$ is a carried tuple $\Delta(e)$ summarizing the authority provenance, interpretation provenance, history / checkpoint basis, mechanism provenance, and randomness / retry lineage behind $e$.
A separation certificate for a pair $(e_1,e_2)$ under a diversity policy $P$ is replayable evidence $\Sep_P(e_1,e_2)$ that the two objects satisfy the policy's non-correlation requirements and therefore count as distinct support.
\end{definition}

\begin{definition}[Support count under policy]
For a finite family $E$ of evidence objects, define
\[
\Supp_P(E)
\]
to be the maximum size of a subset $U\subseteq E$ such that every $e\in U$ locally evaluates to $\Acc$ and every distinct pair in $U$ is $P$-diverse.
A threshold rule ``accept on $t$ supports'' is portable only when it is evaluated using $\Supp_P$, not raw cardinality.
\end{definition}

\begin{definition}[Antichain]
In a poset $(S,\preceq)$, an antichain is a subset $A\subseteq S$ such that no two distinct elements of $A$ are comparable.
\end{definition}

\begin{proposition}[Monotone publication duties admit antichain monitoring (informal)]
Suppose policy-declared publication obligations are monotone over $(S,\preceq)$, and replay carries a sampled antichain together with a domination claim saying that every realized element whose publication matters is covered by at least one sampled antichain element under that same order.
If monitors verify the required inclusion witnesses on the sampled antichain and the carried domination claim is sound,
then any missing required publication yields a missing witness on at least one sampled antichain element.
So antichain sampling is a selection-surface optimization over already-carried publication obligations rather than a family-local acceptance rule.
\end{proposition}

\begin{proof}[Proof sketch]
Monotonicity says that once the sampled antichain covers the realized set in the carried order, every publication duty that matters to replay is witnessed through at least one sampled antichain element.
So if some required publication is missing, the missing obligation cannot stay hidden below the sampled frontier: it must surface as a missing witness for one of the sampled elements that was claimed to cover it.
The monitoring shortcut therefore changes only which carried obligations are checked directly, not the underlying portability rule for those obligations.
\end{proof}

\begin{theorem}[Portable plurality requires disclosed search and disclosed separation]
Assume replay is fail-closed.
Suppose an acceptance decision depends on a selection surface $\Sel$ and a policy $P$ that either
(a) accepts from a menu only after some realized branch path is admissible, or
(b) accepts only when at least $t$ carried supports count distinctly under $P$.
If a bundle carries neither the realized lineage nor a valid branch-invariance certificate for the policy-relevant branch choice,
or if it counts multiple supports without replayable dependency separation, then portable replay may not return $\Acc$ on the strength of that hidden plurality.
In particular, hidden retries, hidden menu search, hidden stopping behavior, and relabeled copies of one unresolved dependency class cannot turn a non-accept into a portable accept.
\end{theorem}

\begin{proof}
Portable replay must determine two things: whether the carried acceptance corresponds to a policy-admissible branch path,
and whether the supports that are supposed to strengthen assurance really count as distinct under policy.
If the realized branch path is hidden and no branch-invariance certificate is carried, replay cannot distinguish
``this outcome would have held regardless'' from ``the producer searched until one branch happened to pass.''
If multiple supports are counted but no carried dependency separation justifies treating them as distinct,
replay cannot distinguish ``several fault domains agreed'' from ``one hidden dependency was repackaged several times.''
Under fail-closed replay, either unresolved ambiguity is non-authorizing.
So hidden search and silent double-counting cannot create a new portable acceptance.
\end{proof}

\begin{proposition}[Optional stopping and retry families are first-class selection obligations]
Suppose evidence accumulates over time and a producer may emit a decision after any prefix, while retries share some mechanism provenance or latent retry lineage.
Then the stopping rule and realized stopping event are part of the carried lineage, the consumed search opportunities must appear in the attempt ledger or search-budget certificate,
and the retry family counts as more than one support only when replay can also see why the attempts are policy-diverse or carries a valid separation certificate.
\end{proposition}

\begin{proof}
Changing the stopping event changes which evidence prefix is judged, so optional stopping is a latent branch choice unless the hidden prefixes are proven outcome-irrelevant.
Likewise, retries that share unresolved dependencies remain one support class until replayable separation evidence says otherwise,
and hidden search opportunities remain unresolved unless the carried accounting shows which of them were used or how many could have been used.
A portable decision that depends on hidden stopping behavior, hidden search budget, or hidden independence therefore fails the same fail-closed plurality check.
\end{proof}

\begin{definition}[Stopping surface]
A stopping surface for sequential evidence is the carried subtuple
\[
\Sigma_{\mathrm{stop}}=(\beta,\varphi,\varrho,\Pi)
\]
of the selection surface recording the budget on admissible prefixes or redraws, the stopping / escalation rule, any randomness or resampling policy that affects those prefixes, and the interpretation profile under which the emitted statistic is read.
\end{definition}

\begin{definition}[Stopping-safe accumulation witness]
A stopping-safe accumulation witness for a sequential claim is carried evidence naming the update rule, threshold rule, and scope identity for a sequential statistic family $(A_t)_{t\ge 0}$ such that replay can reconstruct the same prefix values from the carried evidence stream.
Typical examples are named e-processes, supermartingale scores, or other sequential monitors whose validity is claimed to survive the disclosed stopping rule.
\end{definition}

\begin{proposition}[Portable sequential acceptance needs a carried stopping surface]
Suppose a producer may emit a decision after any admissible evidence prefix.
If the bundle carries neither
(i) the realized stopping event together with the stopping surface, nor
(ii) a branch-invariance certificate showing that every policy-admissible stopping event yields the same authorizing outcome,
then portable replay cannot turn the observed prefix into a new $\Acc$.
If acceptance further relies on a special sequential statistic such as an e-process,
then replay also needs a stopping-safe accumulation witness naming that statistic family under the same stopping surface.
\end{proposition}

\begin{proof}
Without the stopping surface and realized stopping event, replay cannot distinguish
``the disclosed prefix was predetermined or policy-safe'' from ``the producer watched the stream until one prefix happened to cross the line.''
That is hidden search over admissible prefixes.
A branch-invariance certificate can discharge the ambiguity by proving the choice of stopping event outcome-irrelevant.
Otherwise fail-closed replay must block authorization.
If the claimed robustness comes from a special sequential statistic, replay must additionally know which statistic family and threshold rule were in force;
otherwise the producer can narrate an ordinary searched prefix as if it had always been produced by a stopping-safe monitor.
\end{proof}

\paragraph{Conservative role escalation is bridge-owned.}
When a selection or outcome claim relies on a bridge-owned conservative replay-facing role comparison on one fixed scope, the kernel cites the bridge's replay-facing role-package discipline together with the carried conservative-comparison tuple required there.

\section{Outcome interfaces, diagnostic non-accept states, and applicable counterevidence}
The kernel uses $\{\Acc,\Rej,\Unk\}$ as its coarse replay interface, but local evaluators may expose a finer blocked state such as $\Pend$ when evidence is incomplete for explicitly temporal or procedural reasons.
That extra vocabulary does \emph{not} create a second way to authorize reuse.
It only diagnoses why authorization is still blocked.
The same outcome-facing section also owns the rule for replayable counterevidence:
revocations, rollback triggers, superseding ratchets, emergency rejects, stop-ship notices, or debt-limit blockers do not create a second evidence layer; they determine whether an otherwise supportive bundle is still authorizing and, if later discharged, why portability may resume.

\begin{definition}[Outcome interface]
An \emph{outcome interface} is a finite set $\mathcal{O}$ together with an authorization map
\[
\Auth : \mathcal{O}\to\{0,1\}.
\]
In EvidenceVault there is a distinguished authorizing state $\Acc$ with $\Auth(\Acc)=1$.
Every other state is non-authorizing.
\end{definition}

\begin{definition}[Conservative refinement]
Let $(\mathcal{O}_1,\Auth_1)$ and $(\mathcal{O}_2,\Auth_2)$ be outcome interfaces.
A map $F:\mathcal{O}_1\to\mathcal{O}_2$ is a \emph{conservative refinement} if
\[
\Auth_2(F(o))=\Auth_1(o)
\qquad\text{for every }o\in\mathcal{O}_1.
\]
So the authorizing boundary is preserved exactly even if the non-authorizing side is renamed or subdivided.
\end{definition}

Two common local alphabets in the archive are
\[
\mathcal{O}_{\mathsf{kernel}}=\{\Acc,\Rej,\Unk\},
\qquad
\mathcal{O}_{\mathsf{timed}}=\{\Acc,\Rej,\Pend\}.
\]
In both cases only $\Acc$ authorizes portable reuse.
The difference is diagnostic: $\Unk$ emphasizes missing or hidden prerequisites, while $\Pend$ emphasizes closure that is still incomplete but might later become decidable.

\begin{proposition}[Refinement invariance of authorization safety]
Let $T$ be any transformation on evidence states and let $\Eval$ be an evaluator into an outcome interface $(\mathcal{O}_1,\Auth_1)$.
If
\[
\Auth_1(\Eval(T(E))) \le \Auth_1(\Eval(E))
\quad\text{for all evidence states }E,
\]
so $T$ cannot manufacture new authorizations, then the same is true after any conservative refinement $F:\mathcal{O}_1\to\mathcal{O}_2$:
\[
\Auth_2(F(\Eval(T(E)))) \le \Auth_2(F(\Eval(E))).
\]
\end{proposition}

\begin{proof}
A conservative refinement preserves the authorization bit pointwise.
So any downgrade theorem proved at the authorization boundary transfers unchanged to a more diagnostic local outcome alphabet.
\end{proof}

\begin{definition}[Preclusive object]
A \emph{preclusive object} is carried evidence
\[
  b=(s,r,p,h,\Pi,\sigma)
\]
where $s$ is the scope to which the blocker applies, $r$ is the blocker kind, $p$ is the blocking payload or trigger,
$h$ is any required anchored-history context, $\Pi$ is the interpretation/replay profile under which the blocker is read,
and $\sigma$ is the required authentication and authority evidence.
A preclusive object is portable only when those carried fields replay under the same governed replay context as the supportive evidence.
\end{definition}

\begin{definition}[Applicability and discharge]
A preclusive object $b$ is \emph{applicable} to an acceptance instance $I$ when its carried scope, policy context, validity conditions, and anchor view all replay against the same judged object as the supportive evidence.
Discharge evidence $\Dis(b)$ is carried evidence showing, under policy, that $b$ expired by its own terms, was superseded or revoked on the same replayable history, imposed compensating conditions now discharged by receipts, or was overridden by a stronger authorized object in the same scope.
If applicability or discharge is itself uncheckable, fail-closed replay treats the path as non-authorizing.
\end{definition}

\begin{theorem}[Applicable blockers must dominate until discharged]
\label{thm:blockers-dominate}
Assume replay is fail-closed.
If a bundle contains an applicable preclusive object $b$ for an acceptance instance $I$ and no carried discharge evidence accepted by policy,
then portable replay may not return $\Acc$ for $I$.
Deleting $b$, hiding its scope, or omitting the evidence needed to determine whether it applies cannot legitimately turn a non-authorizing state into portable acceptance.
\end{theorem}

\begin{proof}
By applicability, $b$ lies within the same replayable scope, meaning profile, time window, authority chain, and anchor view as the supportive evidence for $I$.
Policy therefore requires either a valid discharge for $b$ or a non-authorizing outcome.
Accepting anyway would amount to treating a required negative premise as nonexistent.
That violates fail-closed replay.
Omitting $b$ or omitting the evidence needed to determine whether $b$ applies does not help: unresolved applicability is itself non-authorizing.
So portable $\Acc$ is impossible until a valid discharge or override is carried.
\end{proof}

\begin{proposition}[Older support does not survive supersession by default]
\label{prop:supersession-no-default-survival}
Suppose a supportive receipt $r$ was portable under policy state $P_t$, and the carried bundle also contains an applicable later blocker $b$ saying that $P_t$ is superseded, revoked, rolled back, or temporarily barred for the relevant scope.
Then $r$ does not remain portable merely because it once authorized reuse; portability now additionally depends on discharge evidence for $b$ or on replay under the superseding policy state.
\end{proposition}

\begin{proof}
The kernel authorizes reuse at the current replay state, not at the historical moment when $r$ was first accepted.
Once $b$ is applicable, $r$ can contribute support only together with carried evidence that the blocker no longer governs the judged scope.
Without that discharge, the outcome stays on the non-authorizing side of the interface.
\end{proof}


\section{Receipt algebra after the replay and selection surfaces are fixed}
Once replay has fixed the admissible omission boundary and the carried selection / stopping surface, receipts need a conservative algebra.
This section fixes that algebra for already-admissible replay objects: conjunction, weakening, thresholds, sequential accumulation, and a replay-facing normal form for composed receipt bundles.

\begin{definition}[Receipt]
A receipt is a tuple
\[
R=(\mathsf{claim},\mathsf{refs},\mathsf{meta},\sigma),
\]
where $\mathsf{refs}$ contains the relevant carried digests (policies, configs, proofs, logs, manifests, checkpoints, or similar replay coordinates), and $\sigma$ authenticates
$\Canon(\mathsf{claim},\mathsf{refs},\mathsf{meta})$ under a public key or keyless identity.
\end{definition}

\begin{remark}[Family-local signed artifacts are not a second receipt species]
Families may carry signed checker outputs, timestamp attestations, monitor statements, or similar local artifacts as ordinary bundle members or payload fields.
The bridge's inline same-scope carrier discipline keeps those signed artifacts inline unless a family explicitly admits them as replay objects here.
They become \emph{kernel receipts} only in that explicit case; otherwise they remain just carried evidence that a family carrier may reference, not a second receipt notion competing with the kernel's definition.
\end{remark}

We write $\mathsf{val}(R)\in\{\Acc,\Rej,\Pend\}$ for the local verifier's evaluation of $R$.
Here $\Pend$ is only a diagnostic refinement of the non-authorizing frontier: it means ``more carried closure is still required,'' not ``soft accept.''

\paragraph{Conjunctive composition.}
When a decision depends on multiple independent receipts --- for example, ``policy check passed,'' ``proof verified,'' and ``publication obligations met'' ---
the default safe operator is fail-closed conjunction.

\begin{definition}[Conjunctive composition]
For receipts $R_1,R_2$,
\[
\mathsf{val}(R_1 \Combine R_2)=
\begin{cases}
\Acc & \text{if }\mathsf{val}(R_1)=\Acc \text{ and } \mathsf{val}(R_2)=\Acc,\\
\Rej & \text{if }\mathsf{val}(R_1)=\Rej \text{ or } \mathsf{val}(R_2)=\Rej,\\
\Pend & \text{otherwise.}
\end{cases}
\]
\end{definition}

\paragraph{Menus and thresholds.}
Systems often allow disclosed alternative evidentiary paths --- multiple monitors, multiple proof systems, multiple vendors, or multiple audit providers.
Once the carried menu surface, realized choice lineage, stopping behavior, dependency profile, and support-counting rule have been fixed,
a safe operator is threshold support with blocker precedence still enforced by the kernel.

\paragraph{Weakening after the replay surface is fixed.}
Family papers often need to weaken a carried claim after conjunction or threshold composition has already fixed the visible replay surface.
This section fixes the algebraic rule once the weakened object is already an admissible replay object.

\begin{definition}[Weakening]
Suppose receipt claims live in a preorder $(\mathcal{C},\preceq)$ where $c'\preceq c$ means that $c'$ is no stronger than $c$
under the pinned policy/profile.
A weakening of a receipt $R=(c,\mathsf{refs},\mathsf{meta},\sigma)$ is any receipt
\[
\mathsf{wk}(R)=(c',\mathsf{refs}',\mathsf{meta}',\sigma')
\]
such that $c'\preceq c$ and the carried references and metadata still justify that implication.
Weakening can forget strength; it cannot manufacture a stronger portable claim.
\end{definition}

\begin{lemma}[Weakening is conservative]
If replay is fail-closed, then weakening cannot manufacture portable acceptance.
It preserves acceptance only when the weaker claim is independently sufficient under the same carried policy and replay assumptions.
\end{lemma}

\begin{proof}
Weakening moves downward in the claim preorder, so it never adds obligations already absent from the stronger claim.
Under fail-closed replay, forgetting strength cannot turn an unsatisfied prerequisite into a satisfied one.
Hence weakening cannot upgrade a non-authorizing state to portable acceptance.
\end{proof}

\begin{remark}
This section reasons only over already-admissible replay objects.
What is fixed here is the algebra for composing, weakening, thresholding, and accumulating those objects once their replay boundary has already been carried.
\end{remark}


\paragraph{Member-index carriers inherit the kernel algebra.}
Some family-local carriers do not introduce a second decision algebra of their own.
They merely name already-admissible replay objects together with local role or dependency metadata and then hand those members to the fixed kernel evaluation rule.

\begin{proposition}[Member-index carriers do not define a second receipt algebra]
\label{prop:member-index-carrier}
Fix one replayable scope.
Suppose a family-local carrier $M$ contains only
(a) content-addressed references to already-admissible replay objects on that scope and
(b) local role, dependency, or provenance metadata telling replay how those members should later be consulted.
Assume replay of $M$ is defined by first expanding the designated members and only then applying the fixed kernel rule for conjunction, weakening, thresholds, plurality, blocker precedence, or complementary-obligation checks.
Then $M$ does not define a second portable receipt algebra on that scope.
It is only a carrier for member selection and local support metadata, and all portability questions for the downstream combination are discharged by the kernel once those named members are replayable.
\end{proposition}

\begin{proof}
By hypothesis, $M$ contributes no new verifier semantics beyond naming replayable members and the metadata that determines how the fixed kernel rule will consult them.
After expansion, the objects being combined are exactly the already-admissible replay objects on that scope, and the evaluation rule is still the kernel's rule rather than a family-local alternative.
So there is no second algebra to validate beyond checking that the member references and their local metadata are well formed on the inherited replay boundary.
Any change in the downstream acceptance behavior must therefore come from changing the carried member set, its carried role/dependency metadata, or the kernel rule itself.
The first two are ordinary carrier changes, and the last is not family-local.
Hence portability of the downstream combination is governed entirely by the kernel once the named members are replayable.
\end{proof}

\paragraph{Sequential accumulation over a carried stopping surface.}
Receipts often arrive over time: monitors emit periodic attestations, transparency logs grow, and audits land after the first decision.
The portability question --- which prefixes were admissible, which stopping rule was in force, and which stopping event actually fired --- is part of the selection-closure rule above.
Once that stopping surface has been carried, the kernel can still recognize special sequential receipts whose evidentiary value survives adaptive stopping.
One convenient family uses e-values.

\begin{definition}[E-value]
A nonnegative random variable $E$ is an e-value for a null hypothesis $H_0$ if $\mathbb{E}_{H_0}[E]\le 1$.
An anytime-valid test rejects $H_0$ when $E\ge 1/\alpha$.
\end{definition}

\begin{theorem}[Optional-stopping safety (standard)]
If $(E_t)_{t\ge 0}$ is a nonnegative supermartingale under $H_0$ with $E_0=1$, then for any (possibly adaptive) stopping time $\tau$,
\[
\Pr_{H_0}\!\left[\sup_{t\le \tau} E_t \ge \tfrac{1}{\alpha}\right] \le \alpha.
\]
Consequently, using the rule ``reject when $E_t\ge 1/\alpha$'' controls false rejection at level $\alpha$ even under optional stopping.
\end{theorem}

\begin{remark}
The kernel is agnostic to how $E_t$ is produced.
Across EvidenceVault, examples include drift tripwires, equivocation audits, and consistency checks that update incrementally as more log data arrives.
When such evidence is anchored to external histories or signed governance objects, portable composition additionally requires the bridge's governed replay discipline.
\end{remark}

\paragraph{Receipt normal forms.}
For offline replay it is useful to record a normal form for composed receipts:
\begin{itemize}[leftmargin=*]
\item a policy/config digest spine,
\item a multiset of signed leaf receipts,
\item any carried material required by the bridge's governed replay discipline for those leaf receipts,
\item any applicable blocker objects (revocations, emergency rejects, rollback triggers, debt-limit notices) together with their discharge evidence if portability is restored,
\item a deterministic evaluation rule (conjunctive gates, disclosed thresholds, anytime thresholds, and blocker precedence), including any diversity policy used to collapse correlated supports.
\end{itemize}
If that normal form is later published compactly or cited through a local surface, the same carried replay assumptions remain in force;
this section fixes only the evaluation rule over the already-admissible normal form.

\section{Closure vector and portable acceptance criterion}

\paragraph{Identity closure.}
The carried evidence must bind to the exact judged object on the replayable scope.
Without identity closure, nearby-object replay remains possible.

\paragraph{Semantic closure.}
The meaning of acceptance must itself be pinned: verifier semantics, policy digest, schema/canonicalization rules,
and any algebra used to combine carried receipts. A proof is not portable if a downstream verifier can silently change what counts as passing.

\paragraph{Selection closure.}
A result that can depend on hidden retries, hidden menus, or hidden selective disclosure is not yet governance-grade.
Selection closure requires those latent degrees of freedom either to be committed ex ante or to be paid for by explicit lineage, receipts, stopping surfaces, and budgets.

\paragraph{Transparency closure.}
Some claims additionally depend on external consistency facts and governed freshness/publication conditions.
These determine whether a local success can be safely reused elsewhere and for how long.
Portable replay therefore requires one coherent governed replay context for those anchored facts, together with any lineage and dependency evidence needed when multiple witnesses, monitors, retries, or menu outcomes are counted together.

For any replay-bearing object $o$ on a fixed scope, define the closure vector
\[
\Clos(o)=(c_{\mathsf{id}},c_{\mathsf{sem}},c_{\mathsf{sel}},c_{\mathsf{tr}})\in\{0,1\}^4,
\]
where each coordinate records whether the corresponding obligation is discharged by carried evidence.

\begin{proposition}[Monotone replay discipline]
If a replay kernel is fail-closed and only upgrades $\Clos(o)$ when it validates additional evidence, then acceptance is monotone in evidence:
adding valid receipts can change $\Unk$ to $\Acc$ or $\Rej$, but deleting required receipts cannot turn a non-accept into acceptance.
\end{proposition}

\begin{proof}
Each closure bit is certified by explicit evidence. Under fail-closed evaluation, absent evidence cannot satisfy a prerequisite.
Therefore removing evidence cannot create a newly satisfied prerequisite, while adding valid evidence can only enable previously blocked checks.
\end{proof}

\begin{theorem}[Portable acceptance criterion]
Suppose a verifier instantiates the replay kernel above and is sound for the carried certificate semantics.
If $\Clos(o)=(1,1,1,1)$ for the carried evidence bundle, then any downstream verifier using the same pinned policy and transparency context
computes the same acceptance outcome on offline replay.
If any component is $0$, the safe default is non-portability: replay may still be informative, but downstream reuse must be downgraded or refused.
\end{theorem}

\begin{remark}
The theorem is intentionally modest. It does not say the underlying world is good; it says the evidence boundary is stable enough that
``acceptable'' means the same thing to the original producer and to a later auditor.
\end{remark}


\begin{definition}[Normalized row declaration]
Fix one family row on one replayable scope after its admitted carrier family has been declared.
Its spine-side declaration records exactly four surfaces:
\[
(R,K,F,C),
\]
where $R$ is the scope-fixing local form, $K$ is the admitted local carrier family on that scope, $F$ is either $\bot$ or the row's conservative comparison label on that scope, and $C$ is either $\bot$ or the row's citation-role label on that same fixed slice.
The spine fixes only that generic row quadruple.
The current mechanism rows and their instantiated labels are printed once, bridge-side, in the bridge-owned family architecture map.
\end{definition}


\end{document}
